\begin{figure}[t]
\centering
\begin{tikzpicture}[x=1cm,y=1cm,>=Stealth,font=\small]
 \draw[->] (0,0)--(6.4,0) node[right,font=\footnotesize] {$n$};
 \draw[->] (0,0)--(0,3.35) node[above] {$\log_2 d$};
 \draw[dashed,thick] (.3,.17)--(5.8,3.0);
 \draw[thick] (.3,.1)--(5.8,2.32);
 \draw[thick] (.3,.18) .. controls (1.2,.51) and (2.9,.69) .. (5.8,.82);
 \node[anchor=west,align=left,font=\footnotesize] at (6.05,2.99)
   {transcript upper bound\\$m\log_2 B_\tau+O(\log n)=O_\eps(n)$};
 \node[anchor=west,align=left,font=\footnotesize] at (6.05,2.02)
   {incidence separation: $n/3+O(1)$\\$m=\Theta_\eps(n)$, $\tau=\Omega_\eps(1)$};
 \node[anchor=west,align=left,font=\footnotesize] at (6.05,.66)
   {Patel's lower bound: $\Omega(n^{1/3})$\\$m,\tau^{-1}=2^{\widetilde O(n^{2/9})}$};
\end{tikzpicture}
\caption{Growth orders of log dimension at fixed learning and representation errors; constants are suppressed. The incidence bound holds for all three dimension notions and meets the exponential order of the transcript bound at its own query resources. The dashed line uses $K=Q=2^n$, $m=\Theta_\eps(n)$, and fixed tolerance; it is not the upper bound at Patel's resources. Patel's curve denotes an ordinary-dimension lower bound, not an upper bound on his class.}
\label{fig:parameters}
\end{figure}
